% ==============================================================================
\section{Fast CASPaxos}
\label{sec:fast-caspaxos}
% ==============================================================================

Fast CASPaxos extends CASPaxos with shared fast ballots.
The piggybacked prepare becomes an atomic promise for the next ballot, which may be proposer-owned or shared.

\begin{figure*}[t]
\centering
{\footnotesize
\protsetup
\noindent
\begin{minipage}[t]{\protcolwidth}
\begin{protbox}{State}
\protsubhdr{Proposer state \textit{(volatile):}}
\begin{protfields}
\textbf{id} & unique proposer identifier ($\textit{id} > 0$)\\
\textbf{round} & monotonically increasing, initially 0\\
\textbf{prepared} & most recently prepared ballot, initially $(1, 0)$\\
\textbf{cachedVal} & last decided value, initially \texttt{null}\\
\end{protfields}
\protsubhdr{Acceptor state \textit{(persistent):}}
\begin{protfields}
\textbf{promised} & highest ballot promised, initially $(1, 0)$\\
\textbf{accepted} & last accepted ballot, or \texttt{null}\\
\textbf{value} & last accepted value, or \texttt{null}\\
\end{protfields}
\end{protbox}
\vspace{\protboxvgap}
\begin{protbox}{Rules for Proposers}
\protsubhdr{Quorums:}
$q_p = q_c = \lfloor N/2 \rfloor + 1$ (prepare and classic accept),\; $q_f = \lceil 3N/4 \rceil$ (fast accept).
\protsubhdr{Procedure:}
\textit{To execute an operation with update function $f$:}
\begin{protimpl}
\item \textbf{Read.}\; Obtain ballot $b$ and \textit{cachedVal}:
  \begin{protrules}
  \item If \textit{prepared} is not \texttt{null}: let $b = \textit{prepared}$;\; set \textit{prepared} to \texttt{null}.
  \item Else: increment \textit{round}; let $b = (\textit{round}, \textit{id})$; send \textsc{Prepare}($b$) to all acceptors. Wait for $q_p$~promises; set \textit{cachedVal} per Value Recovery below. If fewer than $q_p$ promises: set \textit{round} above highest round from any rejections; retry step 1.
  \end{protrules}
\item \textbf{Modify.}\; let $\textit{newVal} = f(\textit{cachedVal})$.
\item \textbf{Write.}\; Choose \textit{nextBallot}: $(b.r{+}1,\, 0)$ to prepare a fast round;\; $(b.r{+}1,\, \textit{id})$ to prepare a classic round.
\item Send \textsc{Accept}($b$, \textit{newVal}, \textit{nextBallot}) to all acceptors. Wait for $q_f$~responses if $b$ is a fast ballot; $q_c$ if $b$ is a classic ballot.
\item On \textsc{Accept} response:
  \begin{protrules}
  \item Quorum success: $\textit{cachedVal} = \textit{newVal}$;\; if a quorum of successful responses reports $\textit{promised} \ge \textit{nextBallot}$ then set $\textit{prepared} = \textit{nextBallot}$, else set \textit{prepared} to \texttt{null}.
  \item Otherwise (rejection or timeout): set \textit{prepared} to \texttt{null}; set \textit{round} above highest round from any rejections; retry from step 1.
  \end{protrules}
\end{protimpl}
\protsubhdr{Value Recovery:}
\textit{When \textsc{Prepare} returns accepted values from a prepare quorum:}
\begin{protrules}
\item If no response contains an \textit{accepted} ballot: set \textit{cachedVal} to \texttt{null}.
\item If the highest \textit{accepted} ballot is \textbf{classic}: $\textit{cachedVal}$ = the corresponding \textit{value}. Only one value can exist at a classic ballot.
\item If the highest \textit{accepted} ballot is \textbf{fast}: tally responses by \textit{value} at that ballot. Let \texttt{maxCount} be the highest tally and set \textit{cachedVal} to the unique maximum, or any tied maximum.
\end{protrules}
\end{protbox}
\end{minipage}%
\hfill
\begin{minipage}[t]{\protcolwidth}
\begin{protbox}{Ballots}
A ballot is a tuple $(r, p)$ where $r$ is a round number and $p$ is the proposer's \textbf{id}. \textbf{Classic} ($p \neq 0$): only proposer $p$ may use; a classic ballot corresponds to at most one value. \textbf{Fast} ($p = 0$): any proposer may propose.
\par\smallskip\noindent
Ordering: $(r_1, p_1) > (r_2, p_2)$ iff $r_1 > r_2$, or $r_1{=}r_2$ and $p_1 > p_2$.
\end{protbox}
\vspace{\protboxvgap}
\begin{protbox}{Prepare RPC}
\textit{Invoked by proposer to read the current value and acquire exclusive write access.}
\protsubhdr{Arguments:}
\begin{protfields}
\textbf{ballot} & ballot $(r, p)$ for this operation\\
\end{protfields}
\protsubhdr{Results:}
\begin{protfields}
\textbf{success} & \texttt{Promise(accepted, value)}\\
\textbf{failure} & \texttt{Reject(promised)}\\
\end{protfields}
\protsubhdr{Acceptor implementation:}
\begin{protimpl}
\item If $\textit{ballot} < \textit{promised}$, reply \texttt{Reject(promised)}.
\item $\textit{promised} = \textit{ballot}$.
\item Reply \texttt{Promise(accepted, value)}.
\end{protimpl}
\end{protbox}
\vspace{\protboxvgap}
\begin{protbox}{Accept RPC}
\textit{Invoked by proposer to commit a value to the register.}
\protsubhdr{Arguments:}
\begin{protfields}
\textbf{ballot} & ballot to write under\\
\textbf{value} & the value to write\\
\textbf{nextBallot} & optional piggybacked prepare; \texttt{null} skips it\\
\end{protfields}
\protsubhdr{Results:}
\begin{protfields}
\textbf{success} & \texttt{Accepted(promised)}\\
\textbf{failure} & \texttt{Reject(promised)}\\
\end{protfields}
\protsubhdr{Acceptor implementation:}
\begin{protimpl}
\item If $\textit{ballot} < \textit{promised}$, reply \texttt{Reject(promised)}.
\item If $\textit{ballot} = \textit{accepted}$:
  \begin{protrules}
  \item If $\textit{value} = v$: reply \texttt{Accepted(promised)} (state unchanged).
  \item If $\textit{value} \neq v$: reply \texttt{Reject(promised)}.
  \end{protrules}
\item $\textit{accepted} = \textit{ballot}$;\; $\textit{value} = v$.
\item $\textit{promised} = \max(\textit{ballot},\, \textit{nextBallot})$.
\item Reply \texttt{Accepted(promised)}.
\end{protimpl}
\end{protbox}
\end{minipage}%
}% footnotesize
\caption{Fast CASPaxos summary. Relative to CASPaxos (Figure~\ref{fig:caspaxos-summary}), ballots with $p = 0$ are shared fast rounds, fast accepts require $\lceil 3N/4 \rceil$ acceptors, fast ballots use first-write-wins with idempotent duplicates, and Value Recovery chooses the unique maximum or any tied maximum.}
\label{fig:fast-caspaxos-summary}
\end{figure*}


% ------------------------------------------------------------------------------
\subsection{Piggybacked Prepare}
\label{sec:piggybacked-prepare}
% ------------------------------------------------------------------------------

In CASPaxos, a successful \textsc{Accept} may carry a piggybacked prepare for the next ballot.
Fast CASPaxos keeps this mechanism but lets the next ballot be either
\begin{itemize}
    \item $(r+1, p)$: a proposer-owned classic ballot;
    \item $(r+1, 0)$: a shared fast ballot.
\end{itemize}

The piggybacked prepare affects only later rounds: it adds no state and does not change the current accept.

\paragraph{Register initialization.}
Acceptors are initialized with promised ballot $(1, 0)$, so the first operation may attempt a fast round immediately.
If it collides, recovery proceeds through an ordinary classic round.

\paragraph{Partial piggybacked-prepare delivery.}
A proposer may reuse a piggybacked ballot only if a quorum of successful \textsc{Accept} responses confirms that promise.
Otherwise it discards the candidate next ballot and falls back to \textsc{Prepare} on its next operation.

% ------------------------------------------------------------------------------
\subsection{Fast Ballots}
\label{sec:fast-ballots}
% ------------------------------------------------------------------------------

Multiple proposers may concurrently submit different values under the same fast ballot, and different acceptors may accept different values.
Safety requires that at most one value be \emph{chosen}, meaning accepted by a fast quorum.
Fast CASPaxos enforces this by having an acceptor accept the first valid request it receives for a fast ballot and reject later requests for that ballot with a different value.
If concurrent proposals split the acceptors, several values may be accepted, but none is chosen unless one reaches a fast quorum.
Because fast quorums intersect, two different values cannot both be chosen in the same fast ballot; if none reaches a fast quorum, a later classic round recovers a safe value from the highest accepted ballot it sees.

\subsubsection{Idempotent Write for Identical Values}
When multiple proposers submit \emph{identical} values under the same fast ballot---as in almost-everywhere-agreement workloads such as Rapid cluster membership~\cite{suresh_atc18}---strict first-write-wins would cause unnecessary rejections.
We therefore let an acceptor acknowledge an \textsc{Accept} request that exactly matches its current accepted ballot and value.
This does not change state and is safe by Section~\ref{sec:correctness}.

\subsubsection{Value Recovery}
\label{sec:value-recovery}
When a proposer runs \textsc{Prepare}, it must recover a value from the highest accepted ballot reported by the prepare quorum.
The recovery rule is best phrased in terms of the values that \emph{could have been decided} at that highest ballot:

\begin{itemize}
    \item \textbf{No prior value.} If no acceptor in the quorum has accepted any value, the proposer is free to propose any value.
    \item \textbf{Classic ballot.} If the highest accepted ballot is a classic ballot, its value is the only possible decided value at that ballot.
    \item \textbf{Fast ballot.} If the highest accepted ballot is a fast ballot, several values may appear at that ballot. In the fixed-quorum setting, the possible decided values are exactly those tied for maximum support in the prepare responses at that ballot. A unique maximum must be recovered; any tied maximum is safe.
\end{itemize}

\noindent A value chosen in a fast ballot was accepted by $q_f$ acceptors; any later prepare quorum intersects that fast quorum in at least $q_f + q_p - N$ acceptors, so any chosen value appears in the recovery evidence.\footnote{This is the counting-quorum specialisation of the \emph{possibleValues} function from the epochs-by-recovery framework~\cite{howard_tr19}.}
A uniquely chosen value therefore appears strictly more often than any competitor, while an exact tie leaves the prepare quorum unable to distinguish among the tied values.

\subsubsection{Conflict Recovery Walkthrough}
Consider $N = 4$ acceptors with fast quorum $q_f = \lceil 3 \cdot 4 / 4 \rceil = 3$, all prepared for fast ballot $(1, 0)$.
Suppose $P_1$ proposes $v_a$ and $P_2$ proposes $v_b$, and the acceptors split 2/2: $a_1,a_2$ accept $v_a$, $a_3,a_4$ accept $v_b$, and each acceptor rejects the other value under the first-write-wins rule.
No value reaches a fast quorum, so no value is chosen.
$P_1$ then runs \textsc{Prepare}($(2, 1)$) and collects responses from, say, $\{a_1, a_2, a_3\}$, which report $v_a$ twice and $v_b$ once at ballot $(1, 0)$.
Because $v_a$ is the unique maximum at the highest fast ballot seen, $P_1$ must recover $v_a$, send \textsc{Accept}($(2, 1),\, v_a$), and commit after a classic quorum of $q_c = 3$ confirms.
Recovery therefore adds two message delays beyond the failed fast round, for a total of 3~RTT from the client's perspective.
Figure~\ref{fig:collision-recovery} illustrates this scenario.

\begin{figure}[t]
\centering
\resizebox{\columnwidth}{!}{%
\begin{tikzpicture}[x=1.12cm, y=1.20cm, font=\footnotesize, >=Stealth]
\def\xEnd{8.8}
\def\yPa{3.00}
\def\yPb{2.40}
\def\yA{1.80}
\def\yB{1.20}
\def\yC{0.60}
\def\yD{0.0}

\node[stproc] at (0,\yPa) {$P_1$};
\node[stproc] at (0,\yPb) {$P_2$};
\node[stproc] at (0,\yA) {$a_1$};
\node[stproc] at (0,\yB) {$a_2$};
\node[stproc] at (0,\yC) {$a_3$};
\node[stproc] at (0,\yD) {$a_4$};

\draw[stlane] (0.25,\yPa) -- (\xEnd,\yPa);
\draw[stlane] (0.25,\yPb) -- (\xEnd,\yPb);
\draw[stlane] (0.25,\yA) -- (\xEnd,\yA);
\draw[stlane] (0.25,\yB) -- (\xEnd,\yB);
\draw[stlane] (0.25,\yC) -- (\xEnd,\yC);
\draw[stlane] (0.25,\yD) -- (\xEnd,\yD);

\node[stsection] at (1.60,3.70) {\textsf{1}};
\node[streqbox] at (1.20,3.12) {\textsc{Accept}$((1{,}0),\, v_a)$};
\node[streqbox] at (1.20,2.52) {\textsc{Accept}$((1{,}0),\, v_b)$};
\draw[streq] (0.85,\yPa) -- (1.55,\yA);
\draw[streq] (0.85,\yPa) -- (1.70,\yB);
\draw[streq] (1.00,\yPb) -- (1.70,\yC);
\draw[streq] (1.00,\yPb) -- (1.85,\yD);
\draw[strsp] (1.55,\yA) -- (2.65,\yPa);
\draw[strsp] (1.70,\yB) -- (2.65,\yPa);
\draw[strsp] (1.70,\yC) -- (2.65,\yPb);
\draw[strsp] (1.85,\yD) -- (2.65,\yPb);
\node[stwarnbox] at (2.45,1.45) {split 2/2};

\draw[dotted, gray] (3.10,-0.16) -- (3.10,3.24);

\node[stsection] at (4.55,3.70) {\textsf{2}};
\node[streqbox] at (4.10,3.12) {\textsc{Prepare}$((2{,}1))$};
\draw[streq] (3.75,\yPa) -- (4.40,\yA);
\draw[streq] (3.75,\yPa) -- (4.55,\yB);
\draw[streq] (3.75,\yPa) -- (4.70,\yC);
\draw[streq] (3.75,\yPa) -- (4.85,\yD);
\draw[strsp] (4.40,\yA) -- (5.75,\yPa);
\draw[strsp] (4.55,\yB) -- (5.75,\yPa);
\draw[strsp] (4.70,\yC) -- (5.75,\yPa);
\node[stnotebox] at (5.45,2.52) {$v_a \times 2,\; v_b \times 1$};

\draw[dotted, gray] (6.15,-0.16) -- (6.15,3.24);

\node[stsection] at (7.45,3.70) {\textsf{3}};
\node[streqbox] at (7.00,3.12) {%
    \textsc{Accept}$((2{,}1),\, v_a)$\\[-1pt]
    \textcolor{blue}{prepare $(3{,}0)$}%
};
\draw[streq] (6.65,\yPa) -- (7.30,\yA);
\draw[streq] (6.65,\yPa) -- (7.45,\yB);
\draw[streq] (6.65,\yPa) -- (7.60,\yC);
\draw[streq] (6.65,\yPa) -- (7.75,\yD);
\draw[strsp] (7.30,\yA) -- (8.55,\yPa);
\draw[strsp] (7.60,\yC) -- (8.55,\yPa);
\node[stcommitbox] at (8.48,\yPa-0.16) {commit $v_a$};

\draw[staxis] (0.25,-0.42) -- node[below, font=\scriptsize] {time} (\xEnd,-0.42);
\end{tikzpicture}%
}% end resizebox
\caption{Fast-round collision and classic recovery. $P_1$ and $P_2$ split the acceptors 2/2 under fast ballot~$(1,0)$, so no value reaches $q_f = 3$. $P_1$ then prepares classic ballot~$(2,1)$, recovers $v_a$ as the unique maximum in the prepare responses, commits it, and piggybacks $(3,0)$ to restore a fast next round.}
\label{fig:collision-recovery}
\end{figure}

% ------------------------------------------------------------------------------
\subsection{Learning Decided Values}
\label{sec:learning-committed-values}
% ------------------------------------------------------------------------------

A proposer may send \textsc{Accept} under ballot $b$ only if it knows the register value for $b$.
It can learn that value by running \textsc{Prepare}, by committing and reusing the piggybacked next ballot, or via a best-effort learn notification from another proposer.
Learn traffic is only an optimization: if it is stale or missing, the proposer falls back to \textsc{Prepare}.
Appendix~\ref{sec:python-implementation} sketches one concrete learn-message format.

% ------------------------------------------------------------------------------
